\section{Análise e discussão}
\label{sec:analise}

A análise será organizada em uma descrição da amostra, uma apresentação das transições observadas e uma discussão das estimações. A descrição deverá informar o período coberto, os critérios de inclusão e a disponibilidade das variáveis. As comparações entre grupos ocupacionais deverão explicitar os denominadores, os pesos utilizados e as diferenças entre exposição à inteligência artificial e características efetivamente observadas dos trabalhadores.

A discussão dos resultados deverá considerar a magnitude das associações, sua incerteza e sua coerência com o desenho da pesquisa. Uma diferença entre períodos não será interpretada automaticamente como consequência da inteligência artificial. Será necessário examinar explicações alternativas, alterações na coleta da PNAD Contínua e possíveis mudanças na composição da amostra, mantendo a distinção entre descrição, inferência e interpretação causal.

Este capítulo é provisório e foi preparado para o exercício ia-6.1. Não apresenta resultados numéricos novos. As tabelas e figuras da versão final deverão ser vinculadas às saídas verificadas do projeto.
\subsection{Organização da análise}

A Figura~\ref{fig:fluxo-pesquisa-pnadc} apresenta as etapas previstas para organizar a pesquisa. O esquema descreve o procedimento de trabalho e não representa resultados empíricos.

\begin{figure}[htb]
  \caption{Fluxo de organização da pesquisa}
  \label{fig:fluxo-pesquisa-pnadc}
  \centering
  \fbox{%
    \begin{minipage}{0.85\textwidth}
      \centering
      Preparação e verificação dos insumos

      $\downarrow$

      Correspondência entre classificações ocupacionais

      $\downarrow$

      Construção dos pares de entrevistas da PNAD Contínua

      $\downarrow$

      Cálculo dos indicadores e estimação das associações

      $\downarrow$

      Verificação dos resultados e discussão das limitações
    \end{minipage}%
  }
  \fonte{Elaboração própria para o exercício ia-6.1.}
\end{figure}

A Tabela~\ref{tab:desfechos-pesquisa-pnadc} sintetiza três desfechos considerados na pesquisa. As definições operacionais completas deverão ser apresentadas na metodologia e conferidas no código de construção das variáveis.

\begin{table}[htb]
  \caption{Desfechos considerados na pesquisa}
  \label{tab:desfechos-pesquisa-pnadc}
  \centering
  \begin{tabular}{p{0.28\textwidth} p{0.60\textwidth}}
    \hline
    \textbf{Desfecho} & \textbf{Descrição para organização da análise} \\
    \hline
    Mudança de ocupação &
    Transição entre classificações ocupacionais nas entrevistas consecutivas, conforme o nível de agregação adotado. \\
    \hline
    Saída do emprego &
    Passagem da condição de ocupado para uma condição sem ocupação na entrevista seguinte. \\
    \hline
    Passagem para a informalidade &
    Transição de uma ocupação classificada como formal para uma classificada como informal, conforme a definição do projeto. \\
    \hline
  \end{tabular}
  \fonte{Elaboração própria a partir do desenho documentado da pesquisa.}
\end{table}